% Part: history
% Chapter: biographies
% Section: gerhard-gentzen

\documentclass[../../../include/open-logic-section]{subfiles}

\begin{document}

\olfileid{his}{bio}{gen}

\olsection{ગેરહાર્ડ ગેન્ટ્ઝન}

\olphoto{gentzen-gerhard}{ગેરહાર્ડ ગેન્ટ્ઝન}

ગેરહાર્ડ ગેન્ટ્ઝન મુખ્યત્વે સંરચનાત્મક સાબિતી-સિદ્ધાંતના
સર્જક તરીકે જાણીતા છે; ખાસ કરીને તેમણે સ્વાભાવિક નિગમન
અને સિક્વન્ટ કલનનાં !!{derivation} તંત્રો રચ્યાં.
તેમનો જન્મ 24 નવેમ્બર 1909ના રોજ જર્મનીના
ગ્રાઇફ્સવાલ્ડમાં થયો હતો. તૈયારીની શાળામાં જતાં પહેલાં
તેઓ ત્રણ વર્ષ ઘરે જ ભણ્યા હતા. શાળામાં અભ્યાસની બાબતમાં
તેઓ પોતાના મોટાભાગના સહાધ્યાયીઓથી પાછળ હતા.
છતાં તેઓ તેજસ્વી વિદ્યાર્થી હતા અને ગણિત માટેની
તેમની વિશેષ પ્રતિભા દેખાતી હતી. તેમની રુચિઓ વિવિધ
હતી: દાખલા તરીકે, તેઓ પોતાની માતા માટે કાવ્યો
અને શાળાના નાટ્યમંચ માટે નાટકો પણ લખતા.

ગેન્ટ્ઝને યુનિવર્સિટીનો અભ્યાસ ગ્રાઇફ્સવાલ્ડ
યુનિવર્સિટીમાં શરૂ કર્યો, પરંતુ પછી ગોટિંગેન,
મ્યુનિખ અને બર્લિન ગયા. 1933માં તેમણે ગોટિંગેન
યુનિવર્સિટીમાંથી હર્માન વાઇલના માર્ગદર્શન હેઠળ
ડૉક્ટરેટ મેળવી. (તેમના મોટાભાગના કાર્યનું
માર્ગદર્શન પૉલ બર્નાયસે કર્યું હતું, પરંતુ
નાઝીઓએ તેમને યુનિવર્સિટીમાંથી બરતરફ કર્યા.)
1934માં ગેન્ટ્ઝને ડેવિડ હિલ્બર્ટના સહાયક તરીકે
કામ શરૂ કર્યું. એ જ વર્ષે તેમણે પોતાના
\emph{Untersuchungen \"{u}ber das logische Schlie\ss en I--II
  [તાર્કિક નિગમન અંગેની તપાસ ૧--૨]} નામના
લેખોમાં સિક્વન્ટ કલન અને સ્વાભાવિક નિગમનનાં
!!{derivation} તંત્રો વિકસાવ્યાં. 1936માં
તેમણે પીઆનોના સ્વયંસિદ્ધાંતોની સુસંગતતા સાબિત કરી.

નાઝીઓ સાથે ગેન્ટ્ઝનનો સંબંધ જટિલ હતો.
તેમના માર્ગદર્શક બર્નાયસને જર્મની છોડવાની
ફરજ પાડવામાં આવી હતી; એ જ સમયે ગેન્ટ્ઝન
નાઝીઓના અર્ધલશ્કરી સંગઠન એસએની યુનિવર્સિટી
શાખામાં જોડાયા. બીજા અનેક જર્મનોની જેમ
તેઓ નાઝી પક્ષના સભ્ય પણ હતા. યુદ્ધ દરમિયાન
તેમણે હવાઈ ગુપ્તચર એકમમાં દૂરસંચાર અધિકારી
તરીકે સેવા આપી. જોકે 1942માં માનસિક ભંગાણને
કારણે તેમને ફરજમાંથી મુક્ત કરાયા. ગેન્ટ્ઝન
ખરેખર નાઝી પક્ષ પ્રત્યે વફાદાર હતા કે
શૈક્ષણિક સફળતા સુનિશ્ચિત કરવા પક્ષમાં
જોડાયા હતા, તે સ્પષ્ટ નથી.

1943માં ગેન્ટ્ઝનને પ્રાગની જર્મન યુનિવર્સિટીના
ગણિત સંસ્થાનમાં શૈક્ષણિક પદની દરખાસ્ત થઈ,
જે તેમણે સ્વીકારી. પરંતુ 1945માં પ્રાગના
નાગરિકોએ જર્મન કબજા સામે બળવો કર્યો.
સોવિયેત દળો શહેરમાં આવ્યાં અને યુનિવર્સિટીના
બધા પ્રાધ્યાપકોની ધરપકડ કરી. નાઝી
સંગઠનોમાં તેમના સભ્યપદને કારણે ગેન્ટ્ઝનને
બળજબરીથી મજૂરી કરાવતી શિબિરમાં લઈ જવાયા.
4 ઑગસ્ટ 1945ના રોજ 35 વર્ષની ઉંમરે
કારાકક્ષમાં કુપોષણથી તેમનું અવસાન થયું.

\begin{reading}
ગેન્ટ્ઝનનું સંપૂર્ણ જીવનચરિત્ર \citet{Menzler-Trott2007}માં
જુઓ. નાઝી શાસન હેઠળના ગણિતશાસ્ત્રીઓ વિશેનું
રસપ્રદ પુસ્તક \citet{Segal2014} ગેન્ટ્ઝનના
જીવનનો ટૂંકો ઉલ્લેખ પણ કરે છે.
તાર્કિક નિગમન અંગેના ગેન્ટ્ઝનના લેખો મૂળ
જર્મન ભાષામાં ઉપલબ્ધ છે
\citep{Gentzen1935a,Gentzen1935b}. તેમના લેખોના
અંગ્રેજી અનુવાદો એક જ ગ્રંથમાં
\citet{Gentzen1969} દ્વારા એકત્ર કરવામાં આવ્યા છે;
તેમાં જીવનચરિત્રાત્મક રૂપરેખા પણ છે.
\end{reading}

\end{document}
